\section{帝国终结时的通信、行政与区域节奏}
宣统时期的铁路、议政、军队、报刊、电报与谈判形成一套快速却不均衡的通信网络。本节从行政传递和区域节奏考察帝国终结，避免把北京或武汉的消息时间直接投射为所有省份、县城和家庭的政治时间。
\subsection{1900—1909：区域网络、制度媒介与环境条件}
这些事件的先后可以由账本定位，但每一项影响都须回到具体区域、材料与行动者，不能被压缩为抽象的国家趋势。
\hypertarget{civic:EQ-1909-PROVINCIAL-ASSEMBLIES}{}\paragraph{宣统元年（1909年）}各省咨议局开办，地方精英获得新的公开政治场域。本事件应放在电报、铁路、报刊、军队调动与省级财政的具体网络中理解；信息发出、到达、刊登与转化为行动是不同阶段。 核心取证为《宣统政纪》宣统元年相关条；宫中朱批奏折、内阁大库题本与《清史稿》相关志传。行政文书显示的是机构如何把道路、港口、人口、宗教或案件转换为可处理的类别；没有后续县档或物质证据，不能由分类推定实际生活。账本提示的研究线索为宫廷政治、制度运作与中央—地方信息链研究。事件所处条件为预备立宪需要有限度吸纳地方意见。后续变化应限于选举资格、代表性和实际权力均受限制。来源限制：以宣统元年（1909年）的同期文书为定位起点；官修记录、奏报或外交文本服务特定机构立场，具体日期、规模、地方执行与对手视角须以独立材料复核。
\hypertarget{civic:EQ-1909-RAILWAY-RECOVERY-MOVEMENT}{}\paragraph{宣统元年（1909年）}铁路国有化和路权问题引发各省商绅、公司和官府的争论。本事件应放在电报、铁路、报刊、军队调动与省级财政的具体网络中理解；信息发出、到达、刊登与转化为行动是不同阶段。 核心取证为《宣统政纪》宣统元年相关条；户部奏销、盐政、河工或海关档案。行政文书显示的是机构如何把道路、港口、人口、宗教或案件转换为可处理的类别；没有后续县档或物质证据，不能由分类推定实际生活。账本提示的研究线索为财政账册、市场网络、灾害环境与区域经济研究。事件所处条件为外债、地方集资与中央财政需求相冲突。后续变化应限于“收回路权”诉求包含不同地区和利益群体。来源限制：以宣统元年（1909年）的同期文书为定位起点；官修记录、奏报或外交文本服务特定机构立场，具体日期、规模、地方执行与对手视角须以独立材料复核。
\subsection{1910—1919：区域网络、制度媒介与环境条件}
这些事件的先后可以由账本定位，但每一项影响都须回到具体区域、材料与行动者，不能被压缩为抽象的国家趋势。
\hypertarget{civic:EQ-1910-ZIZHENG-YUAN}{}\paragraph{宣统二年（1910年）}资政院开院，中央层面的立宪咨询机构开始运作。本事件应放在电报、铁路、报刊、军队调动与省级财政的具体网络中理解；信息发出、到达、刊登与转化为行动是不同阶段。 核心取证为《宣统政纪》宣统二年相关条；宫中朱批奏折、内阁大库题本与《清史稿》相关志传。行政文书显示的是机构如何把道路、港口、人口、宗教或案件转换为可处理的类别；没有后续县档或物质证据，不能由分类推定实际生活。账本提示的研究线索为宫廷政治、制度运作与中央—地方信息链研究。事件所处条件为朝廷试图展示宪政进程并控制代表参与。后续变化应限于咨询权与立法权的界限成为持续争议。来源限制：以宣统二年（1910年）的同期文书为定位起点；官修记录、奏报或外交文本服务特定机构立场，具体日期、规模、地方执行与对手视角须以独立材料复核。
\hypertarget{civic:EQ-1910-FAMINE-AND-UNREST}{}\paragraph{宣统二年（1910年）}灾荒、物价和财政压力持续影响晚清地方社会与政治动员。本事件应放在电报、铁路、报刊、军队调动与省级财政的具体网络中理解；信息发出、到达、刊登与转化为行动是不同阶段。 核心取证为《宣统政纪》宣统二年相关条；地方志、督抚奏折及地方档案。编年条目可以为行动建立相对次序，却通常不保存劳动过程、物价变动、建筑材料或非精英的叙述。账本提示的研究线索为地方档案、迁徙、宗教、族群与社会动员研究。事件所处条件为自然环境、市场波动和行政能力交互作用。后续变化应限于地方灾情与政治行动不可简单建立单因果关系。来源限制：以宣统二年（1910年）的同期文书为定位起点；官修记录、奏报或外交文本服务特定机构立场，具体日期、规模、地方执行与对手视角须以独立材料复核。
\hypertarget{civic:EQ-1911-RAILWAY-NATIONALIZATION}{}\paragraph{宣统三年（1911年）}清廷宣布铁路干线国有化并以外债筹资，保路运动迅速扩大。本事件应放在电报、铁路、报刊、军队调动与省级财政的具体网络中理解；信息发出、到达、刊登与转化为行动是不同阶段。 核心取证为《宣统政纪》宣统三年相关条；宫中朱批奏折、内阁大库题本与《清史稿》相关志传。行政文书显示的是机构如何把道路、港口、人口、宗教或案件转换为可处理的类别；没有后续县档或物质证据，不能由分类推定实际生活。账本提示的研究线索为宫廷政治、制度运作与中央—地方信息链研究。事件所处条件为中央财政、外债与地方商办铁路利益冲突。后续变化应限于政策公告、地方抗议和后续镇压应分开记录。来源限制：以宣统三年（1911年）的同期文书为定位起点；官修记录、奏报或外交文本服务特定机构立场，具体日期、规模、地方执行与对手视角须以独立材料复核。
\hypertarget{civic:EQ-1911-SICHUAN-RAILWAY-PROTEST}{}\paragraph{宣统三年（1911年）}四川保路运动升级，赵尔丰处置引发更大政治危机。本事件应放在电报、铁路、报刊、军队调动与省级财政的具体网络中理解；信息发出、到达、刊登与转化为行动是不同阶段。 核心取证为《宣统政纪》宣统三年相关条；地方志、督抚奏折及地方档案。编年条目可以为行动建立相对次序，却通常不保存劳动过程、物价变动、建筑材料或非精英的叙述。账本提示的研究线索为地方档案、迁徙、宗教、族群与社会动员研究。事件所处条件为地方股东、商绅、学生和官府对路权与财政有不同诉求。后续变化应限于参与规模和暴力过程须依当地档案与报刊核验。来源限制：以宣统三年（1911年）的同期文书为定位起点；官修记录、奏报或外交文本服务特定机构立场，具体日期、规模、地方执行与对手视角须以独立材料复核。
\hypertarget{civic:EQ-1911-WUCHANG-UPRISING}{}\paragraph{宣统三年（1911年）}武昌起义爆发，新军与革命团体控制武汉部分地区。本事件应放在电报、铁路、报刊、军队调动与省级财政的具体网络中理解；信息发出、到达、刊登与转化为行动是不同阶段。 核心取证为《宣统政纪》宣统三年相关条；军机处档、战事方略及地方奏报。编年条目可以为行动建立相对次序，却通常不保存劳动过程、物价变动、建筑材料或非精英的叙述。账本提示的研究线索为战争动员、军需、战场暴力与地方社会研究。事件所处条件为新军组织、革命网络和铁路危机相互叠加。后续变化应限于起义的即时参与者、城市控制与后续省份响应需分层叙述。来源限制：以宣统三年（1911年）的同期文书为定位起点；官修记录、奏报或外交文本服务特定机构立场，具体日期、规模、地方执行与对手视角须以独立材料复核。
\hypertarget{civic:EQ-1911-PROVINCIAL-INDEPENDENCE}{}\paragraph{宣统三年（1911年）}多省宣布脱离清廷，帝国行政和财政网络快速分裂。本事件应放在电报、铁路、报刊、军队调动与省级财政的具体网络中理解；信息发出、到达、刊登与转化为行动是不同阶段。 核心取证为《宣统政纪》宣统三年相关条；宫中朱批奏折、内阁大库题本与《清史稿》相关志传。行政文书显示的是机构如何把道路、港口、人口、宗教或案件转换为可处理的类别；没有后续县档或物质证据，不能由分类推定实际生活。账本提示的研究线索为宫廷政治、制度运作与中央—地方信息链研究。事件所处条件为武昌起义、地方军政力量和立宪派选择共同推动变化。后续变化应限于“独立”声明不等于各省形成稳定统一控制。来源限制：以宣统三年（1911年）的同期文书为定位起点；官修记录、奏报或外交文本服务特定机构立场，具体日期、规模、地方执行与对手视角须以独立材料复核。
\hypertarget{civic:EQ-1911-YUAN-SHIKAI-RETURN}{}\paragraph{宣统三年（1911年）}清廷起用袁世凯处理北方军政和议和事务。本事件应放在电报、铁路、报刊、军队调动与省级财政的具体网络中理解；信息发出、到达、刊登与转化为行动是不同阶段。 核心取证为《宣统政纪》宣统三年相关条；宫中朱批奏折、内阁大库题本与《清史稿》相关志传。行政文书显示的是机构如何把道路、港口、人口、宗教或案件转换为可处理的类别；没有后续县档或物质证据，不能由分类推定实际生活。账本提示的研究线索为宫廷政治、制度运作与中央—地方信息链研究。事件所处条件为中央缺乏可依赖的军队与谈判人选。后续变化应限于袁在清廷、北洋与革命派之间的策略塑造退位进程。来源限制：以宣统三年（1911年）的同期文书为定位起点；官修记录、奏报或外交文本服务特定机构立场，具体日期、规模、地方执行与对手视角须以独立材料复核。
\hypertarget{civic:EQ-1912-REPUBLIC-FOUNDING}{}\paragraph{宣统四年（1912年）}中华民国临时政府在南京成立，与清廷退位谈判并行。本事件应放在电报、铁路、报刊、军队调动与省级财政的具体网络中理解；信息发出、到达、刊登与转化为行动是不同阶段。 核心取证为《宣统政纪》宣统四年相关条；宫中朱批奏折、内阁大库题本与《清史稿》相关志传。行政文书显示的是机构如何把道路、港口、人口、宗教或案件转换为可处理的类别；没有后续县档或物质证据，不能由分类推定实际生活。账本提示的研究线索为宫廷政治、制度运作与中央—地方信息链研究。事件所处条件为各省独立与南北议和创造新的政权竞争格局。后续变化应限于临时政府的法理主张与实际控制范围需要区分。来源限制：以宣统四年（1912年）的同期文书为定位起点；官修记录、奏报或外交文本服务特定机构立场，具体日期、规模、地方执行与对手视角须以独立材料复核。
\hypertarget{civic:EQ-1912-QING-ABDICATION}{}\paragraph{宣统四年（1912年）}宣统帝颁布退位诏书，清帝国终结。本事件应放在电报、铁路、报刊、军队调动与省级财政的具体网络中理解；信息发出、到达、刊登与转化为行动是不同阶段。 核心取证为《宣统政纪》宣统四年相关条；宫中朱批奏折、内阁大库题本与《清史稿》相关志传。行政文书显示的是机构如何把道路、港口、人口、宗教或案件转换为可处理的类别；没有后续县档或物质证据，不能由分类推定实际生活。账本提示的研究线索为宫廷政治、制度运作与中央—地方信息链研究。事件所处条件为南北谈判、财政军政压力与皇室待遇安排共同促成退位。后续变化应限于退位文本规定的优待与地方政治现实须分别追踪。来源限制：以宣统四年（1912年）的同期文书为定位起点；官修记录、奏报或外交文本服务特定机构立场，具体日期、规模、地方执行与对手视角须以独立材料复核。
